\chapter*{Introduction Générale}
\addcontentsline{toc}{chapter}{Introduction Générale}

L’ère numérique a profondément transformé le secteur bancaire et financier,
rendant les transactions électroniques plus rapides et le traitement des données plus volumineux.
Cette évolution a également entraîné une augmentation significative des risques de fraude bancaire ~\cite{mckinsey2019},nécessitant le développement de méthodes de détection plus sophistiquées et adaptées aux flux massifs d’informations.

Le présent rapport s’intéresse à la \textbf{détection de fraude bancaire avant et après l’apparition du Big Data}, 
en mettant en évidence les changements technologiques et méthodologiques dans la collecte, le traitement et l’analyse des données. 
L’objectif est de montrer comment les nouvelles approches permettent d’améliorer la précision, la rapidité et l’efficacité des systèmes de détection de fraude.

Le projet a été réalisé au sein de la \textbf{Faculté des Sciences de Tunis (FST)}, 
dans le cadre du module \textbf{Big Data}, sous l’encadrement de \textbf{Mme Manel Zekri}, 
enseignante-chercheure au Département d’Informatique. 
L’objectif principal était de \textbf{comparer les méthodes traditionnelles de détection de fraude avec les approches modernes basées sur le Big Data}, 
tout en concevant un système capable d’identifier les transactions suspectes et de visualiser les résultats de manière interactive et compréhensible.

La période du projet s’étend du \textbf{15/11/2025 au 02/12/2025}, 
au cours de laquelle différentes étapes ont été réalisées : étude du domaine bancaire, 
collecte et traitement des données, conception des modèles de détection, 
et réalisation de visualisations interactives à l’aide d’outils de Business Intelligence.
\clearpage
Le rapport est structuré en cinq parties principales :  

\begin{enumerate}
    \item \textbf{Chapitre 1 : Présentation du projet} — Cadre, problématique, objectifs et choix méthodologiques.
    \item \textbf{Chapitre 2 : Conception du système} — Présentation des données, des acteurs, des processus de détection de fraude et structure générale du système.
    \item \textbf{Chapitre 3 : Étude comparative avant / après Big Data} — Analyse des méthodes traditionnelles et des approches modernes, accompagnée d’un tableau comparatif synthétique.
    \item \textbf{Chapitre 4 : Visualisation des données} — Présentation des tableaux de bord et des indicateurs clés permettant d’illustrer les résultats de la détection de fraude.
    \item \textbf{Conclusion et perspectives} — Synthèse des travaux réalisés et propositions d’amélioration pour de futures études.
\end{enumerate}
